\cleardoublepage
\chapter{PENDAHULUAN}
\label{chap:pendahuluan}

% --- Latar Belakang ---
\section{Latar Belakang}

Pertumbuhan populasi global dan urbanisasi yang pesat telah meningkatkan volume produksi sampah secara signifikan. Menurut World Bank, produksi sampah global diproyeksikan mencapai 3,4 miliar ton pada tahun 2050, meningkat dari 2,01 miliar ton pada tahun 2016 \autocite{kaza2018}.
Peningkatan volume sampah ini menimbulkan tantangan kompleks dalam pengelolaan limbah yang efisien dan berkelanjutan.

Salah satu tantangan utama dalam pengelolaan sampah adalah ketiadaan data terstruktur mengenai pola pembuangan. Pertanyaan seperti kapan terjadi puncak pembuangan, apakah ada perbedaan pola antara hari kerja dan akhir pekan, serta jenis sampah apa yang mendominasi, menjadi sulit dijawab tanpa sistem \textit{monitoring} yang memadai. Metode konvensional seperti pengecekan visual atau pencatatan manual memiliki keterbatasan dalam hal konsistensi, akurasi, dan cakupan waktu, sehingga pengambilan keputusan sering dilakukan tanpa basis data.

Perkembangan teknologi \textit{Internet of Things} (IoT) membuka peluang untuk mengatasi permasalahan tersebut. Sensor yang dipasang pada tempat sampah dapat memantau tingkat kepenuhan secara periodik, menyimpan data ke basis data, dan menghasilkan \textit{insight} mengenai pola pembuangan. Di Indonesia, salah satu contoh perangkat \textit{Smart Waste Bin} yang telah dikembangkan adalah produk dari PT Layanan Cerdas Indonesia (LCI), yang telah diimplementasikan di beberapa kota untuk pemantauan kapasitas kontainer sampah \autocite{LCI2024}. Namun, perangkat \textit{Smart Waste Bin} yang tersedia secara komersial umumnya dirancang untuk skala perkotaan, sehingga konfigurasi perangkat termasuk \textit{firmware} belum tentu sesuai untuk lingkungan yang lebih kecil seperti kampus, di mana dinamika pembuangan sampah berubah dalam hitungan jam dan memerlukan resolusi data yang lebih tinggi.

Berdasarkan permasalahan tersebut, penelitian ini bertujuan untuk mengembangkan sistem \textit{Smart Waste Bin} berbasis IoT yang disesuaikan dengan kebutuhan lingkungan kampus, baik dari sisi perangkat lunak sistem maupun \textit{firmware} perangkat. Sistem yang dibangun kemudian diuji di Gedung Labtek VIII ITB untuk mengumpulkan data dan menganalisis pola pembuangan sampah yang dihasilkan.

% --- Rumusan Masalah ---
\section{Rumusan Masalah}

Berdasarkan latar belakang dan permasalahan di atas, penelitian ini merumuskan pertanyaan penelitian sebagai berikut:

\begin{enumerate}
    \item Bagaimana rancangan dan implementasi sistem \textit{Smart Waste Bin} berbasis IoT yang mampu mengumpulkan data pembuangan sampah secara periodik dan terstruktur?

    \item Pola pembuangan sampah apa yang dapat diidentifikasi berdasarkan data yang dikumpulkan oleh sistem, dan seberapa valid pola tersebut terhadap observasi lapangan?
\end{enumerate}

% --- Tujuan ---
\section{Tujuan}

Penelitian ini bertujuan untuk merancang dan mengimplementasikan sistem \textit{Smart Waste Bin} berbasis \textit{Internet of Things} untuk memahami pola pembuangan sampah di Gedung Labtek VIII ITB. Secara detil, tujuan penelitian ini adalah:

\begin{enumerate}
    \item Merancang arsitektur dan mengimplementasikan sistem \textit{Smart Waste Bin} berbasis IoT yang mampu mengumpulkan data pembuangan sampah secara periodik dan terstruktur.

    \item Mengidentifikasi pola pembuangan sampah berdasarkan data yang dikumpulkan oleh sistem dan memvalidasinya melalui observasi lapangan.
\end{enumerate}

% --- Batasan Masalah ---
\section{Batasan Masalah}

Penelitian ini memiliki batasan-batasan berikut:

\begin{enumerate}
    \item Pengujian lapangan dari sistem yang dibangun dilakukan di Lantai 4 area depan Laboratorium IoT, Gedung Labtek VIII Institut Teknologi Bandung.

    \item Sistem menggunakan 3 unit perangkat \textit{Smart Waste Bin} yang masing-masing ditempatkan berdasarkan jenis sampah: organik, anorganik, dan residu.

    \item Perangkat IoT yang digunakan merupakan modifikasi dari perangkat \textit{Smart Waste Bin} pihak ketiga yang disesuaikan untuk kebutuhan lingkungan kampus.

    \item Interval pengumpulan data ditetapkan $\pm$2,5 jam berdasarkan pertimbangan kapasitas daya baterai perangkat.

    \item Klasifikasi jenis sampah mengacu pada penempatan fisik tempat sampah (organik, anorganik, residu), bukan melalui sensor klasifikasi otomatis.

    \item Analisis pola pembuangan sampah mencakup pola temporal (distribusi jam dan hari), tren volume, serta komposisi jenis sampah berdasarkan data yang dikumpulkan sistem.
\end{enumerate}

% --- Metodologi ---
\section{Metodologi}

Pelaksanaan tugas akhir ini didukung oleh studi literatur dan menggunakan metodologi \textit{Design Science Research} (DSR) untuk pengembangan sistem. Studi literatur dilakukan untuk membangun fondasi teoritis dan memahami teknologi yang relevan, sedangkan DSR digunakan sebagai kerangka kerja pengembangan dan evaluasi artefak yang berfokus pada pemenuhan kebutuhan (\textit{requirement}) yang telah ditetapkan.

\subsection{Studi Literatur}

Studi literatur dilakukan untuk mengumpulkan informasi yang relevan tentang topik yang diangkat, mencakup teori, konsep, teknologi, dan penelitian terdahulu yang telah dilakukan oleh peneliti lain. Studi literatur ini menjadi fondasi dalam pengembangan sistem \textit{Smart Waste Bin} berbasis \textit{Internet of Things}.

Aktivitas yang dilakukan pada tahap studi literatur meliputi:

\begin{enumerate}
    \item \textbf{Pencarian Literatur}

    Melakukan pencarian literatur dari berbagai sumber akademik dan ilmiah, seperti jurnal internasional, konferensi, artikel ilmiah, buku, dan dokumentasi teknis. Pencarian dilakukan menggunakan basis data akademik seperti IEEE \textit{Xplore}, \textit{Google Scholar}, \textit{ScienceDirect}, dan \textit{Springer}. Kata kunci yang digunakan dalam pencarian mencakup: ``\textit{smart waste bin}'', ``IoT \textit{waste management}'', ``\textit{waste pattern analysis}'', ``\textit{smart waste monitoring}'', ``\textit{sensor-based waste management}'', dan kombinasi kata kunci lainnya yang relevan.

    \item \textbf{Pengelompokan Literatur}

    Literatur yang diperoleh dikelompokkan berdasarkan kategori-kategori seperti konsep dan teori dasar tentang \textit{Internet of Things} (IoT), teknologi sensor dan perangkat keras untuk \textit{monitoring} sampah, arsitektur sistem dan integrasi IoT, penelitian terdahulu tentang \textit{smart waste bin} dan pengelolaan sampah berbasis teknologi, dan metode analisis data temporal.

    \item \textbf{Penapisan dan Evaluasi Literatur}

    Literatur yang telah dikumpulkan disaring berdasarkan relevansi, kualitas publikasi, dan kebaruan informasi.

    \item \textbf{Sintesis dan Dokumentasi}

    Hasil studi literatur disintesis untuk mengidentifikasi \textit{gap} penelitian, \textit{best practices}, dan pendekatan yang dapat diadopsi atau dikembangkan lebih lanjut. Seluruh proses studi literatur didokumentasikan dan hasil sintesisnya akan dijelaskan secara sistematis di Bab II Studi Literatur.
\end{enumerate}

\subsection{Metodologi \textit{Design Science Research}}

Pengembangan sistem \textit{Smart Waste Bin} menggunakan metodologi \textit{Design Science Research} (DSR) sebagaimana dirumuskan oleh Peffers dkk. \autocite{peffers2007}. Metodologi ini dipilih karena penelitian ini berorientasi pada pembangunan dan evaluasi sebuah artefak, yaitu sistem \textit{Smart Waste Bin}, bukan pada riset perilaku pengguna. DSR menyediakan kerangka kerja yang sistematis dan teruji untuk penelitian yang menghasilkan artefak teknologi, dengan validasi yang difokuskan pada pemenuhan kebutuhan (\textit{requirement}) yang telah ditetapkan. DSR mencakup enam fase utama, yaitu identifikasi masalah dan motivasi, definisi tujuan solusi, perancangan dan pengembangan, demonstrasi, evaluasi, dan komunikasi. Berikut adalah penjelasan untuk setiap fase yang diterapkan secara spesifik pada penelitian ini:

\begin{enumerate}
    \item \textbf{Identifikasi Masalah dan Motivasi}

    Fase ini bertujuan untuk mengidentifikasi permasalahan yang melatarbelakangi penelitian beserta motivasinya. Aktivitas pada fase ini meliputi:
    \begin{enumerate}[a.]
        \item Mengidentifikasi permasalahan pengelolaan sampah yang masih bergantung pada penjadwalan tetap tanpa mempertimbangkan kondisi aktual kepenuhan tempat sampah, serta ketiadaan data dan informasi kondisi sampah yang terstruktur.
        \item Melakukan pengamatan langsung terhadap kondisi sistem pembuangan yang ada untuk memperkuat pemahaman terhadap permasalahan.
    \end{enumerate}

    \item \textbf{Definisi Tujuan Solusi}

    Fase ini bertujuan untuk merumuskan tujuan solusi berdasarkan permasalahan yang telah diidentifikasi, yang dinyatakan dalam bentuk kebutuhan sistem. Aktivitas pada fase ini meliputi:
    \begin{enumerate}[a.]
        \item Mengolah hasil identifikasi masalah untuk merumuskan \textit{gap} dan kebutuhan akan sistem pengumpulan data yang terstruktur.
        \item Merumuskan spesifikasi kebutuhan sistem, yang mencakup pendefinisian Kebutuhan Fungsional (pengumpulan data, penyimpanan, dan visualisasi pola) serta Kebutuhan Nonfungsional (performa dan reliabilitas). Spesifikasi kebutuhan inilah yang menjadi acuan utama dalam perancangan dan evaluasi.
    \end{enumerate}

    \item \textbf{Perancangan dan Pengembangan}

    Fase ini bertujuan untuk merancang dan membangun artefak sistem \textit{Smart Waste Bin} berdasarkan kebutuhan yang telah ditetapkan. Aktivitas pada fase ini meliputi:
    \begin{enumerate}[a.]
        \item Mengeksplorasi berbagai alternatif arsitektur dan solusi teknologi, serta memilih alternatif terbaik menggunakan metode pengambilan keputusan yang terukur berdasarkan kriteria kelayakan teknis, kesesuaian tujuan, waktu, dan biaya.
        \item Merancang dan mengimplementasikan adaptasi pada sisi \textit{firmware} perangkat keras bawaan untuk menyesuaikannya dengan karakteristik lingkungan kampus dan mengoptimalkan keandalan operasional.
        \item Mengembangkan sistem penerimaan, transmisi, dan penyimpanan data pembuangan sampah sesuai dengan arsitektur yang terpilih.
        \item Mengintegrasikan seluruh komponen perangkat keras dan perangkat lunak sehingga membentuk satu kesatuan sistem yang dapat beroperasi secara otomatis.
    \end{enumerate}

    \item \textbf{Demonstrasi}

    Fase ini bertujuan untuk menunjukkan bahwa artefak yang telah dibangun dapat digunakan untuk menyelesaikan permasalahan, melalui penerapan di lapangan. Aktivitas pada fase ini meliputi:
    \begin{enumerate}[a.]
        \item Menempatkan unit \textit{smart waste bin} pada lokasi yang telah ditentukan untuk mengumpulkan data pembuangan sampah dalam periode tertentu.
        \item Melakukan \textit{monitoring} terhadap operasional sistem selama periode pengumpulan data untuk memastikan sistem merekam dan mentransmisikan data dengan konsisten.
    \end{enumerate}

    \item \textbf{Evaluasi}

    Fase ini bertujuan untuk mengukur seberapa baik artefak memenuhi kebutuhan yang telah ditetapkan pada fase definisi tujuan. Aktivitas pada fase ini meliputi:
    \begin{enumerate}[a.]
        \item Mengevaluasi pemenuhan Kebutuhan Fungsional dan Nonfungsional sistem berdasarkan data hasil demonstrasi.
        \item Menganalisis data yang terkumpul menggunakan analisis statistik deskriptif dan analisis temporal untuk memahami pola pembuangan sampah di lokasi penelitian.
        \item Mengevaluasi efektivitas adaptasi \textit{firmware} yang telah dikembangkan terhadap keandalan dan kualitas data sistem.
    \end{enumerate}

    \item \textbf{Komunikasi}

    Fase ini bertujuan untuk mengomunikasikan permasalahan, artefak yang dihasilkan, dan hasil evaluasinya. Aktivitas pada fase ini berupa penyusunan dokumentasi penelitian secara sistematis dalam bentuk laporan tugas akhir, yang mencakup seluruh proses mulai dari identifikasi masalah hingga kesimpulan.
    \vfill
\end{enumerate}

% --- Sistematika Penulisan ---
\section{Sistematika Penulisan}

Laporan tugas akhir ini disusun dengan sistematika penulisan sebagai berikut:

\begin{enumerate}
    \item Bab I Pendahuluan: Berisi latar belakang, rumusan masalah, tujuan, batasan masalah, metodologi, dan sistematika penulisan laporan tugas akhir.

    \item Bab II Studi Literatur: Berisi kajian pustaka yang relevan dengan topik tugas akhir, mencakup teori dan konsep dasar \textit{Internet of Things}, teknologi sensor dan perangkat keras untuk \textit{monitoring} sampah, arsitektur sistem IoT, serta penelitian terdahulu yang berkaitan dengan \textit{smart waste bin} dan pengelolaan sampah berbasis teknologi.

    \item Bab III Analisis: Berisi penjelasan tentang analisis kebutuhan pengguna dan sistem, termasuk pendefinisian kebutuhan fungsional dan nonfungsional sistem \textit{Smart Waste Bin}.

    \item Bab IV Perancangan: Berisi penjelasan tentang perancangan solusi yang diusulkan, mencakup arsitektur sistem IoT 4-\textit{layer}, perancangan perangkat keras, perangkat lunak, dan \textit{firmware}.

    \item Bab V Implementasi: Berisi penjelasan tentang proses implementasi sistem, termasuk pengembangan \textit{firmware}, integrasi komponen, dan pengembangan \textit{dashboard} visualisasi data.

    \item Bab VI Evaluasi: Berisi analisis dan evaluasi terhadap sistem yang diimplementasikan, mencakup pengujian kebutuhan fungsional dan nonfungsional, analisis pola pembuangan sampah, serta evaluasi peningkatan \textit{firmware}.

    \item Bab VII Kesimpulan dan Saran: Berisi kesimpulan dari hasil pelaksanaan tugas akhir dan saran untuk pengembangan lebih lanjut.

    \item Daftar Pustaka: Berisi daftar referensi yang digunakan dalam penyusunan laporan tugas akhir.

    \item Lampiran: Berisi dokumen-dokumen pendukung seperti data lengkap hasil pengujian dan dokumentasi lapangan.
\end{enumerate}